\subsection{The D01 obstruction made exact (new in v6)}
\label{sec:HI}

Let \(\nu=2-\alpha^{-1}\in(0,1)\).
In the flipped weight-two channel, the CR--F difference of radial
heat kernels is
\[
\Delta k(s;r,r')=k^-_\nu-k^+_\nu
=\frac{1}{2s}e^{-(r^2+r'^2)/(4s)}\,
\frac{2\sin\pi\nu}{\pi}\,K_\nu\!\left(\frac{rr'}{2s}\right).
\]
The D01 derivative tower carries this channel into \(S=3\) through
the raising operator
\[
X=\partial_r-\frac{\nu}{r},
\]
the radial part of the helicity-raising derivative on the
\(J_{-\nu}\) branch.
The term that \((\mathrm H_{I,\mathrm{D01}})\) must control is the
composite diagonal
\[
\mathcal R(s;r)=\bigl[X_rX_{r'}\Delta k(s;r,r')\bigr]_{r'=r}.
\]

\begin{theorem}[Exact obstruction and Mellin value for D01]
\label{thm:HI}
\leavevmode
\begin{enumerate}[label=\textup{(\roman*)}]
\item \emph{Obstruction.}
As \(r\to0\), \(\mathcal R(s;r)=c_*\,r^{-2\nu-2}+O(r^{-2\nu})\),
with \(c_*\neq0\) for every \(0<\nu<1\).
Hence \(\int_0^\epsilon r\,\mathcal R\,\dd r\) diverges.
No \(L^2\) intertwiner and no trace-class statement holds for the
composite. The unrestricted \(\HI\) fails, exactly as the v4 record
states.

\item \emph{Mellin value.}
The Mellin transform
\(M(z)=\int_0^\infty r^{1+z}\mathcal R(s;r)\,\dd r\)
converges for \(\Re z>2\nu\).
It continues meromorphically, with poles contained in
\((2\nu-2\NN_0)\cup(-2\nu-2\NN_0)\).
So it is regular at \(z=0\) for \(0<\nu<1\), and
\[
M(0)=0\qquad\text{identically in }\nu\in(0,1).
\]
\end{enumerate}
\end{theorem}

\begin{proof}
Put \(s=1/2\), so that \(x=r^2\); by scale invariance
\(M(0)\) scales as \(s^{-1}\).
Then
\[
\mathcal R=A\,e^{-x}\Bigl[xK-2xK'+K'+xK''
-2\nu(K'-K)+\nu^2x^{-1}K\Bigr],
\qquad K=K_\nu(x),
\]
where \(A=\sin(\pi\nu)/(\pi s)\).
Use \(K'=-\tfrac12(K_{\nu-1}+K_{\nu+1})\) and
\(K''=\tfrac14(K_{\nu-2}+2K_\nu+K_{\nu+2})\).
Each term then has the form
\(\int_0^\infty x^{\mu-1}e^{-x}K_m(x)\,\dd x\) with
\[
\int_0^\infty x^{\mu-1}e^{-x}K_m(x)\,\dd x
=\frac{\sqrt\pi\,\Gamma(\mu+m)\Gamma(\mu-m)}
{2^\mu\,\Gamma(\mu+\tfrac12)},
\qquad \Re\mu>\abs{\Re m}.
\]
This is the \(a\to1\) limit of DLMF~10.43.22 \cite{DLMF}.
\textsc{[exec]} It is checked by quadrature in its strip to
\(10^{-16}\) (\loc{logs/v6/hi_d01_exact.json}).
The right side is meromorphic in \(\mu\), which gives the stated
continuation and poles.
Leading-order expansion of \(K_\nu\) gives (i).

At \(z=0\) the bracket is a sum of ten gamma ratios with
\(\mu\in\{0,1,2\}\) and \(m\in\{\nu,\nu\pm1,\nu\pm2\}\).
\textsc{[exec]} Applying \(\Gamma(w+1)=w\Gamma(w)\) reduces it
to \(0\) as a symbolic function of \(\nu\)
(\loc{code/hi_d01_exact.py}).
The sum at \(z=10^{-30}\) is \(O(10^{-30})\) for seven values of
\(\nu\).

An independent route uses direct quadrature with the
\(r^{-2\nu-2}\) and \(r^{-2\nu}\) heads subtracted and continued
analytically (\loc{code/hi_d01.py}, 60 digits).
It gives \(s\,M(0)\in[-1.8\times10^{-21},\,3.7\times10^{-32}]\)
at \(\alpha=0.6,0.7,0.75,0.8,0.9\), and the value is independent
of \(s\).
The residual grows like \(\epsilon^{4-2\nu}\) for the cutoff
\(\epsilon=10^{-9}\), consistent with zero.
\end{proof}

\begin{corollary}[Status of \((\mathrm H_{I,\mathrm{D01}})\)]
\label{cor:HI}
Interpret the apex coefficient of the forbidden D01 composite by
homogeneous (Mellin) continuation.
This is the unique scale-covariant extension of the
non-integrable diagonal here, since there is no resonance for
\(0<\nu<1\).
With that reading the composite contributes \(0\), and the D01
constant-profile entries of Table~\ref{tab:hhC} need no correction
from it.

The hypothesis is therefore \emph{upgraded} in two respects: its
failure as an \(L^2\) statement is now exact, and the Mellin value
is exactly zero.
One thing is not shown: that the XVI-P prescription requires
homogeneous continuation of this composite, as opposed to some
other extension of the divergent apex integral.
That is a prescription statement, part of (R3)--(R5), and it is
recorded as \CONDITIONAL.
\end{corollary}
